The project uses only public, zero-budget data: raw files come
from the manual Stack Exchange Data Explorer (SEDE) browser
interface or the public Stack Exchange API v2.3 (no paid quotas,
BigQuery, or private data). Raw files are immutable; derived
tables are regenerated from Python under \texttt{src/}, and the
full pipeline is documented in \texttt{docs/replication\_notes.md}.

\subsection{Stack Overflow tag-week-question-type panel}

The core product is a panel at the (tag, week, question type)
level, 2020-01-01 through 2024-12-31, covering \textbf{100
pre-treatment top tags} (fixed from the pre-ChatGPT period),
\textbf{261 calendar weeks}, and \textbf{7 derived question
types} generated client-side from title-keyword patterns:
short\_code, long\_code, how\_to, debugging\_simple,
other\_conceptual, version\_environment\_specific, and
advanced\_architecture. The first five are coded
\emph{substitutable} ($\text{Sub}_k=1$) on the hypothesis that
current LLMs answer them well. Of the two non-substitutable
categories, version\_environment\_specific is the genuinely
clean control (LLMs cannot reliably reproduce user-specific
stack traces, environment configurations, or repository state),
while advanced\_architecture is a declared \emph{boundary
category}---partial LLM coverage on its conceptual side,
organisation-specific context on its applied side. We declare
this boundary status from the outset and report two baselines
(\S\ref{sec:results}): the binary specification coding
advanced\_architecture as non-substitutable ($-0.108$) and the
boundary-excluded specification using
version\_environment\_specific as the sole control ($-0.140$),
the more defensible reference.

The master panel has 164{,}351 cells, built from 8{,}404{,}411
raw question-tag pairs across 5{,}000{,}258 unique questions.
Observed annual question counts fall from 2.7 million (2020) to
0.5 million (2024)---a stylized fact that motivates the
identification design. Until 2023-05 the panel was built from
parameterized SEDE SQL pasted into the browser; from 2023-06 the
same questions are pulled from the public API (a free registered
key, identical output schema), validated byte-for-byte against
SEDE on the May 2023 overlap (3{,}364 rows per side, identical
across every metadata column).

\paragraph{2025 extension (full panel).}
To test whether the association persists beyond the one-off
2023--2024 platform events (\S\ref{ssec:id_threats}), we
collected calendar-2025 questions for \emph{all 100 tags} from
the same public API. The collection is complete and audited
(100/100 tags, 0 failed windows, 0 page-cap flags, $87{,}027$
unique 2025 questions); the combined 2020--2025 panel has
$186{,}188$ cells over 313 weeks and reproduces the original
panel exactly before 2024-12-30, supporting a full
re-estimation, not just a persistence check.

\subsection{Treatment intensity: AI-answerability}
\label{ssec:data_ai_answerability}

The treatment is a tag-level \emph{proxy} built from
pre-treatment features alone (to avoid mechanical contamination):
pre-ChatGPT accepted-answer rate, historical question frequency,
short-code share, how-to share, and tag maturity. Four variants
are available ($z$-score, PCA, quartile rank, and a conservative
structural variant excluding text-derived features); results are
reported across all four. Because a historical proxy is open to
the objection that it captures popularity rather than
answerability, we run the design on \emph{two co-primary}
measures and require the result under both: the structural
historical index ($-0.108$, $p=0.021$) and an embedding-based
score that uses \emph{no} historical features---built only from
the semantic similarity of question titles to fixed
LLM-answerable versus LLM-resistant exemplars under
\texttt{all-MiniLM-L6-v2} (free, no API cost)---which gives
$-0.119$ ($p<0.001$). The embedding measure shares none of the
proxy's inputs, so the popularity objection cannot explain its
estimate; that the two agree is why we read the gradient as
answerability-driven. The two converge moderately (Pearson
$r\in[0.41,0.50]$ across variants) without being
interchangeable, and the proxy also correlates with an objective
\emph{revealed-answerability} score on 30{,}241 sampled questions
(tag-level Pearson $r=0.563$). Full validation tables (embedding,
partial-regression, revealed-answerability, and AUC by stratum,
which shows the question-level signal concentrated in the
lowest-answerability quartile) are in the replication package; we
are explicit that the instrument is a tag-level intensity, not a
question-level classifier.

\subsection{Outcomes and governance}

The main outcome is $\log(1+Q_{t,w,k})$; we also report
$\log(1+\text{unique users})$, fractional tag counts, and PPML on
raw counts. A zero-budget exploratory GH Archive panel is
documented in the replication notes but not reported here, as it
is uninformative for the present margin (too few language
clusters and censored first-seen contributors). Stack Overflow content is CC BY-SA 4.0; the
project distributes only derived statistics, never underlying
user content.
